\subsection{Commutative stellar algebras of endomorphisms of a Hilbert space}

DEFINITION 3. --- Let A be an algebra over \(\mathbf{C}\) and E a complex topological vector space. Let π be a representation of A in E. An element x of E is said to be a cyclic vector for π if the set of elements π(a)x for a ∈ A is total in E.

Let \(u \in \mathcal{L}(\mathrm{E})\) be an endomorphism of E. We say that \(x \in \mathrm{E}\) is a cyclic vector for \(u\) if it is a cyclic vector for the identity representation of the stellar subalgebra generated by \(u\) in \(\mathcal{L}(\mathrm{E})\) .

PROPOSITION 8. --- Let E be a complex Hilbert space. Let A be a commutative unital stellar subalgebra of \(\mathcal{L}(\mathrm{E})\) . Let \(x\) be an element of E. Set \(\mathrm{E}_x = \overline{\mathrm{A} \cdot x}\) and let us denote \(\mu_x\) the spectral measure of \(x\) relative to A.

There exists a unique isometric isomorphism \(\theta_x\) of \(\mathrm{L}^2 (\mathsf{X}(\mathrm{A}),\mu_x)\) onto \(\mathrm{E}_x\) such that \(\theta_{x}(f) = \mathcal{H}_{\mathrm{A}}(f)(x)\) for every function \(f\in \mathcal{C}(\mathsf{X}(\mathrm{A}))\) . For every element \(a\in \mathrm{A}\) , the space \(\mathrm{E}_x\) is stable under \(a\) and if we denote by \(a_{x}\) the endomorphism of \(\mathrm{E}_x\) deduced from a by passing to subspaces, then we have \(a_{x}\circ \theta_{x} = \theta_{x}\circ m_{\mathcal{G}_{\mathrm{A}}(a)}\) .

Let \(\widetilde{\theta}_x\) be the linear map from \(\mathcal{C}(\mathrm{X(A)})\) into \(\mathrm{E}_x\) defined by

\[
\widetilde {\theta} _ {x} (f) = \mathcal {H} _ {\mathrm{A}} (f) (x).
\]

For any function \(f \in \mathcal{C}(\mathsf{X}(\mathsf{A}))\) , we have

\[
\| \widetilde {\theta} _ {x} (f) \| ^ {2} = \langle \mathcal {H} _ {\mathrm{A}} (f) x | \mathcal {H} _ {\mathrm{A}} (f) x \rangle = \langle x | \mathcal {H} _ {\mathrm{A}} (| f | ^ {2}) x \rangle = \int_ {\mathsf {X} (\mathrm{A})} | f | ^ {2} d \mu_ {x}.
\]

Since \(\mathcal{C}(\mathrm{X(A)})\) is dense in \(\mathcal{L}^2 (\mathrm{X(A)},\mu_x)\) , there exists a unique isometric linear map from \(\mathrm{L}^2 (\mathrm{X(A)},\mu_x)\) into \(\mathrm{E}_x\) which extends \(\widetilde{\theta}_x\) ; it is denoted by \(\theta_{x}\) . The map \(\theta_{x}\) is surjective since its image is closed (Lemma 8 of I, p. 107) and contains \(\mathrm{A}\cdot x\) . It is therefore an isometric isomorphism.

Let \(a \in \mathrm{A}\) . Denote by \(g = \mathcal{G}_{\mathrm{A}}(a) \in \mathcal{C}(\mathsf{X}(\mathrm{A}))\) . The space \(\mathrm{E}_x\) is stable under \(a\) . For \(f \in \mathcal{C}(\mathsf{X}(\mathrm{A}))\) , we then have

\[
\widetilde {\theta} _ {x} (m _ {g} (f)) = \mathcal {H} _ {\mathrm{A}} (\mathcal {G} _ {\mathrm{A}} (a) f) x = (a \circ \mathcal {H} _ {\mathrm{A}} (f)) x = a (\widetilde {\theta} _ {x} (f)),
\]

and it follows that \(\theta_x \circ m_g = a_x \circ \theta_x\) .

COROLLARY. — Let E be a complex Hilbert space. Let A be a commutative unital stellar subalgebra of \(\mathcal{L}(\mathrm{E})\) admitting a cyclic vector \(x\) . Let \(\mu_x\) be the spectral measure of \(x\) with respect to A. There exists a unique isometric isomorphism

\[
\theta_ {x} \colon \mathrm{L} ^ {2} (\mathsf {X} (\mathrm{A}), \mu_ {x}) \to \mathrm{E}
\]

such that \(\theta_x(f) = \mathcal{H}_{\mathrm{A}}(f)\) for every function \(f \in \mathcal{C}(\mathsf{X}(\mathrm{A}))\) . For every a in A, we have \(a \circ \theta_x = \theta_x \circ m_{\mathcal{G}_{\mathrm{A}}(a)}\) .

Indeed, we then have \(E_{x}=E\) and \(a_{x}=a\) for all \(a\in A\) .

Let E be a complex Hilbert space. Let \(x \in E\) and let A be a commutative unital stellar subalgebra of \(\mathcal{L}(\mathrm{E})\) . Denote by \(E_{x}\) the closed subspace \(\overline{A \cdot x}\) of E. For every \(a \in A\) , we then have \(a(\mathrm{E}_{x}) \subset \mathrm{E}_{x}\) . Denote by \(A_{x}\) the commutative unital stellar subalgebra of \(\mathcal{L}(\mathrm{E}_{x})\) consisting of the endomorphisms of \(E_{x}\) deduced from the elements of A by passing to subspaces. The vector x is a cyclic vector for \(A_{x} \subset \mathcal{L}(\mathrm{E}_{x})\) .

PROPOSITION 9. --- There exists a subset C of E such that E is the Hilbert sum of the spaces \(E_{x}\) for \(x \in C\) . If E is countably generated, the set C is countable.

Let \(\mathcal{O}\) the set of subsets C of E such that the subspaces \(\mathrm{E}_x\) for \(x\in \mathrm{C}\) are pairwise orthogonal. The set \(\mathcal{O}\) , ordered by inclusion, is of finite character (E, III, p. 34, def. 2) since C belongs to \(\mathcal{O}\) if and only if the sets formed of two elements of C belong to \(\mathcal{O}\) . By E, III, p. 35, th. 1, there exists a maximal element C of \(\mathcal{O}\) .

Let F be the closed subspace of E generated by the subspaces \(E_{x}\) for \(x \in C\) . It suffices to show that \(F^{\circ}\) is reduced to 0 to complete the proof of the proposition. Let y be an element of \(F^{\circ}\) . For every \(x \in C\) and all a and b in A, we have \(\langle a(y) | b(x) \rangle = \langle y | a^{*}b(x) \rangle = 0\) since \(a^{*}b(x)\) belongs to \(E_{x}\) , hence in F. Since the elements \(a(y)\) (resp. \(b(x)\) ) generate a dense subspace of \(E_{y}\) (resp. \(E_{x}\) ), we therefore have \(E_{y} \subset E_{x}^{\circ}\) . Since C is maximal in O, this means that \(E_{y}\) is reduced to 0, hence that y = 0.

Suppose that E is countably generated. Since \(E_{x}\) is nonzero for all \(x \in C\) non-zero, any set C such that E is the Hilbert sum of the spaces \(E_{x}\) for \(x \in C\) is countable.

THEOREM 1. --- Let A be a commutative unital stellar subalgebra of \(\mathcal{L}(\mathrm{E})\) . There exists a locally compact topological space X, a positive measure \(\mu\) on X, an isometric isomorphism \(\theta\) of \(\mathrm{L}^2 (\mathrm{X},\mu)\) onto E and an isometric morphism of stellar algebras \(\pi\) from A into \(\mathcal{C}_b(\mathrm{X})\) such that for every \(a\in \mathrm{A}\) , one has

\[
a \circ \theta = \theta \circ m _ {\pi (a)}.
\]

By prop. 9, there exists a subset C of E such that E is the Hilbert sum of the subspaces \(E_{x}\) for \(x \in C\) . Let X be the locally compact topological space \(\mathsf{X}(\mathsf{A}) \times \mathsf{C}\) , where C is equipped with the discrete topology. There exists a unique measure \(\mu\) on X such that \(\mu|(\mathsf{X}(\mathsf{A}) \times \{x\})\) is identified with the spectral measure \(\mu_{x}\) of x for every \(x \in C\) (INT, III, p. 65, § 2, n°1, prop. 1); this measure is positive (cf. loc. cit.). One then has a Hilbert sum decomposition

\[
\mathrm{L} ^ {2} (\mathrm{X}, \mu) = \bigoplus_ {x \in \mathrm{C}} \mathrm{L} ^ {2} (\mathrm{X} (\mathrm{A}), \mu_ {x})
\]

(prop. 1 of IV, p. 179, c) applied to the sets \(\mathsf{X}(\mathsf{A}) \times \{x\}\) for \(x \in \mathrm{C}\) ). There therefore exists a unique isometry \(\theta: \mathrm{L}^2(\mathrm{X}, \mu) \to \mathrm{E}\) which coincides with \(\theta_x: \mathrm{L}^2(\mathsf{X}(\mathsf{A}), \mu_x) \to \mathrm{E}_x\) on \(\mathrm{L}^2(\mathsf{X}(\mathsf{A}), \mu_x)\) for all \(x \in \mathrm{C}\) .

Let \(p \colon \mathrm{X} \to \mathrm{X}(\mathrm{A})\) be the canonical projection; it is surjective, hence the linear map \(p^* \colon \mathcal{C}_b(\mathrm{X}(\mathrm{A})) \to \mathcal{C}(\mathrm{X})\) defined by \(f \mapsto f \circ p\) is injective. Denote \(\pi = p^* \circ \mathcal{G}_{\mathrm{A}}\) ; this is an injective morphism of stellar algebras from A into \(\mathcal{C}_b(\mathrm{X})\) ; it is therefore isometric (prop. 9 of I, p. 112).

Let \(a \in \mathrm{A}\) . For every \(x \in \mathrm{C}\) , we have \(a_x \circ \theta_x = \theta_x \circ m_{\mathcal{G}_A(a)}\) . The continuous linear maps \(a \circ \theta\) and \(\theta \circ m_{\pi(a)}\) therefore coincide on the subspaces \(\mathrm{L}^2(\mathrm{X}(\mathrm{A}), \mu_x)\) , and are consequently equal.

COROLLARY (Spectral Theorem). --- Let E be a complex Hilbert space. Let \(u \in \mathcal{L}(\mathrm{E})\) be a normal endomorphism of E. There exist a locally compact topological space X, a positive measure \(\mu\) on X, an isometric isomorphism \(\theta\) of \(\mathrm{L}^{2}(\mathrm{X},\mu)\) onto E and a continuous bounded function g on X such that \(u = \theta \circ m_{g} \circ \theta^{-1}\) .

If u admits a cyclic vector x, one may take \(X = \mathrm{Sp}(u)\) and for g the identity function of X.

It suffices to apply Theorem 1 (resp. the corollary of Proposition 8) to the stellar subalgebra A of \(\mathcal{L}(\mathrm{E})\) generated by \(u\) , and to set \(g = \pi(u)\) .

This statement reduces every question concerning a single normal endomorphism of a Hilbert space to a similar question for a multiplication operator, which often simplifies its study (cf. for example Exercise 19 of IV, p. 325).

